% ============================================
% LUC-01
% Gerado de molambudos.md — build_miolo.py
% ============================================
\fragtitle{luc}{L}{LUC}{01}{O Convite}
\markboth{LUC-01 --- O Convite}{}

\noindent \textit{Inverno de 2026 --- Belo Horizonte}

\noindent \textbf{Inverno de 2026 --- Belo Horizonte}

\vspace{0.5em}\noindent\rule{\textwidth}{0.5pt}\vspace{0.5em}

\noindent O e-mail chegou numa terça-feira, 8:47 da manhã.

\noindent Dra. Lúcia Mendes estava no consultório, com uma xícara de café esfriando na mesa e três laudos para revisar. Ela clicou no remetente: \textit{Arquivo Público Mineiro --- Setor de Documentos Históricos}.

\noindent \textit{\textquotedblleft{}Prezada Dra. Lúcia,}

\noindent \textit{Conforme conversamos em janeiro, localizamos o prontuário completo do Paciente 1.260 (Hospital Colônia de Barbacena) e os anexos mencionados: relatórios do Dr. Heitor Oliveira, laudo de necropsia do IML, e o diário do paciente.}

\noindent \textit{Os documentos estão disponíveis para consulta em nosso arquivo. Considerando seu interesse acadêmico no caso, disponibilizamos uma sala para estudo.}

\noindent \textit{Atenciosamente,}
\textit{Dr. Ronaldo Torres --- Diretor do Arquivo Público Mineiro\textquotedblright{}}

\noindent Lúcia leu o e-mail três vezes.

\noindent Ela tinha solicitado o acesso ao arquivo em janeiro --- uma pesquisa de rotina para um artigo sobre transtornos dissociativos em populações institucionalizadas. O caso do Colônia era famoso na psiquiatria forense brasileira: sessenta mil mortos em oitenta anos, pacientes enterrados vivos, diagnósticos fabricados.

\noindent Mas o Paciente 1.260 era diferente.

\noindent Ela tinha encontrado uma menção ao caso numa nota de rodapé da tese de doutorado de alguém --- uma linha: \textit{\textquotedblleft{}Paciente 1.260, Joaquim Antônio Correia, internado em 1917, faleceu em 1979. Diário encontrado junto ao corpo na exumação. Conteúdo não analisado.\textquotedblright{}}

\noindent Conteúdo não analisado.

\noindent Aquilo a intrigou. O que um paciente do Colônia poderia ter escrito que \textit{ninguém} analisou em quase cinquenta anos?

\noindent \textit{\textquotedblleft{}Você vai se arrepender\textquotedblright{}}, disse a voz na cabeça dela. Não era uma voz sobrenatural. Era a voz da experiência. Quinze anos de psicologia forense tinham lhe ensinado que alguns arquivos deveriam permanecer fechados.

\noindent Ela pegou o telefone.

\noindent \textit{\textquotedblleft{}Arquivo Público Mineiro, bom dia.\textquotedblright{}}

\noindent \textit{\textquotedblleft{}Dra. Lúcia Mendes. Confirme meu agendamento para consulta ao prontuário 1.260.\textquotedblright{}}

\noindent \textit{\textquotedblleft{}Um momento...\textquotedblright{}\textit{ Teclas. }\textquotedblleft{}Dra. Lúcia, o senhor Ronaldo deixou uma observação aqui. Ele pede que a senhora ligue para ele antes de vir. Disse que é... urgente.\textquotedblright{}}

\noindent \textit{\textquotedblleft{}Urgente?\textquotedblright{}}

\noindent \textit{\textquotedblleft{}Sim. Ele disse: 'Avise a Dra. Lúcia que o diário não está mais no arquivo.'\textquotedblright{}}

\noindent Lúcia segurou o telefone com mais força.

\noindent \textit{\textquotedblleft{}Como assim, não está mais?\textquotedblright{}}

\noindent \textit{\textquotedblleft{}Ele explica pessoalmente. Mas pediu para avisar: o diário foi retirado do arquivo em 1979. Por um médico. O Dr. Heitor Oliveira. O mesmo que assina os relatórios anexos ao prontuário.\textquotedblright{}}

\noindent \textit{\textquotedblleft{}Retirado quando?\textquotedblright{}}

\noindent \textit{\textquotedblleft{}Dois dias antes de desaparecer.\textquotedblright{}}

\noindent O café na mesa de Lúcia tinha esfriado completamente.

\noindent \textit{\textquotedblleft{}Estou indo para aí agora.\textquotedblright{}}

\noindent \textit{\textquotedblleft{}Dra. Lúcia...\textquotedblright{}\textit{ A voz do atendente hesitou. }\textquotedblleft{}O senhor Ronaldo disse outra coisa. Disse: 'Avise que o diário foi devolvido semana passada. Por uma pessoa que não se identificou. Devolvido pessoalmente. Na mesa da recepção. Envolto num pano preto.'\textquotedblright{}}

\noindent \textit{\textquotedblleft{}...E?\textquotedblright{}}

\noindent \textit{\textquotedblleft{}Ela disse: 'Avise a Dra. Lúcia que o diário está esperando ela.'\textquotedblright{}}

\noindent Lúcia desligou.

\noindent Olhou para o café frio.
Olhou para os laudos.
Olhou para a hora.

\noindent 8:52.

\noindent Cinco minutos depois do e-mail, ela já sabia que algo estava errado.

\noindent Ela só não sabia \textit{o quanto}.

\vspace{0.5em}
\noindent\textit{\textbf{↪ Rotas:} \rota{LUC-02} \rotasep \rota{DOC-08}}
